%%%%%%%%%%%%%%%%%%%%%%%%%%%%%%%%%%%%%%%%%%%%%%%%%%%%%%%%%%%%%%%%%%%%%%%%%%%%%%%
% SECTION 6: Language Processing
%%%%%%%%%%%%%%%%%%%%%%%%%%%%%%%%%%%%%%%%%%%%%%%%%%%%%%%%%%%%%%%%%%%%%%%%%%%%%%%
\section{Language Processing}

% The Limits of N-grams
\begin{frame}{The Old Approach: N-grams}
    \textbf{How NLP worked before deep learning:}
    
    \vspace{1em}
    \begin{itemize}
        \item Count frequencies of word sequences of length $N$
        \item Predict next word by frequency alone
        \item Possible sequences scale as $V^N$ where $V$ = vocabulary size
    \end{itemize}
    
    \vspace{1em}
    \textbf{The problems:}
    \begin{itemize}
        \item Hopelessly unscalable
        \item No semantic understanding
        \item ``fast car'' and ``quick car'' are completely unrelated
        \item Fundamentally counting, not understanding
    \end{itemize}
\end{frame}

% Word Embeddings - The Insight
\begin{frame}{Word Embeddings: The Insight}
    \textbf{Bengio et al.'s approach:}
    
    \vspace{1em}
    \begin{enumerate}
        \item Represent each word as \textbf{one-hot vector}
        
        (one component = 1, rest = 0)
        
        \item First layer learns to map these into \textbf{dense embeddings}
        \item Semantically similar words end up geometrically close
    \end{enumerate}
    
    \vspace{1em}
    % TODO: Add embedding space visualization
    
    \textbf{This structure was not programmed — it emerged:}
    \begin{quote}
        \small``When trained to predict the next word in a news story, the learned word vectors for Tuesday and Wednesday are very similar, as are the word vectors for Sweden and Norway.''
    \end{quote}
\end{frame}

% Word Embeddings - Why It's Surprising
\begin{frame}{Word Embeddings: Why It's Surprising}
    \textbf{It works on:}
    \begin{itemize}
        \item Small, clean datasets (Hinton's 1986 family trees)
        \item Large, messy real-world corpora
    \end{itemize}
    
    \vspace{1em}
    \textbf{This was not the expected result at the time}
    
    \vspace{1em}
    \textbf{Foreshadows:}
    \begin{itemize}
        \item Word2Vec (2013)
        \item GloVe (2014)
        \item Every modern language model embedding
    \end{itemize}
\end{frame}

% RNNs - The Problem
\begin{frame}{Recurrent Neural Networks: The Problem}
    \textbf{Feedforward networks for language:}
    \begin{itemize}
        \item Stateless — each input processed independently
        \item No memory of prior context
        \item Cannot model sequence, grammar, long-range dependencies
    \end{itemize}
    
    \vspace{1em}
    \textbf{Language is inherently sequential}
    
    \vspace{0.5em}
    ``The cat sat on the \_\_\_'' requires knowing what came before
\end{frame}

% RNN Structure
\begin{frame}{RNN: Unrolled Structure}
    % TODO: Add RNN unrolled diagram showing:
    % - Input sequence x₁, x₂, x₃, ...
    % - Hidden states h₁, h₂, h₃, ...
    % - Outputs y₁, y₂, y₃, ...
    % - Arrows showing h_{t-1} → h_t connections
    % - IMPORTANT: Label that weights W_hh, W_xh, W_hy are SHARED across all time steps
    
    \begin{center}
        \textit{[Diagram: RNN unrolled through time]}
    \end{center}
    
    \vspace{1em}
    \textbf{Key insight:} The \textbf{same weights} are used at every time step
    $$h_t = \sigma(W_{hh} h_{t-1} + W_{xh} x_t + b)$$
    
    \vspace{0.5em}
    This weight sharing is both the power and the problem...
\end{frame}

% RNNs - The Solution
\begin{frame}{RNNs: What They Enable}
    \textbf{Key idea:} Pass a hidden state forward through the sequence
    
    \vspace{1em}
    \begin{itemize}
        \item Network retains implicit memory of previous inputs
        \item Enables sentence-level context
        \item Sequence-to-sequence tasks become possible
    \end{itemize}
    
    \vspace{1em}
    \textbf{Example: Machine translation}
    \begin{itemize}
        \item Translate English → French
        \item Learns implicit grammatical and semantic structure
        \item Not word-by-word lookup
    \end{itemize}
\end{frame}

% The Vanishing Gradient Problem
\begin{frame}{The Vanishing Gradient Problem}
    \textbf{Backpropagating through time:}
    
    \vspace{0.5em}
    $$\frac{\partial L}{\partial h_1} = \frac{\partial L}{\partial h_T} \cdot \prod_{t=2}^{T} \frac{\partial h_t}{\partial h_{t-1}}$$
    
    \vspace{1em}
    \begin{itemize}
        \item Each $\frac{\partial h_t}{\partial h_{t-1}}$ involves $W_{hh}$ and $\sigma'$
        \item \textbf{Same weights multiplied $T$ times}
        \item If $\|W_{hh}\| < 1$: gradients \textbf{vanish} exponentially
        \item If $\|W_{hh}\| > 1$: gradients \textbf{explode}
    \end{itemize}
    
    \vspace{1em}
    \textbf{Result:} Long-range dependencies cannot be learned
    
    \vspace{0.5em}
    \small\textit{``The cat, which sat on the mat that belonged to my neighbor, \_\_\_''}
\end{frame}

% LSTM Structure
\begin{frame}{LSTM: Structure}
    \textbf{Hochreiter \& Schmidhuber, 1997}
    
    % TODO: Add LSTM cell diagram showing:
    % - Cell state C_t (the "conveyor belt" running through)
    % - Forget gate f_t
    % - Input gate i_t  
    % - Output gate o_t
    % - Candidate values C̃_t
    % - How gates control information flow
    
    \begin{center}
        \textit{[Diagram: LSTM cell with gates labeled]}
    \end{center}
    
    \vspace{1em}
    \textbf{Key innovation:} Separate \textbf{cell state} $C_t$ from hidden state $h_t$
    
    Cell state flows through with minimal transformation — gradients preserved
\end{frame}

% LSTM Gates
\begin{frame}{LSTM: The Gates}
    \textbf{Forget gate:} What to discard from cell state
    $$f_t = \sigma(W_f \cdot [h_{t-1}, x_t] + b_f)$$
    
    \vspace{0.5em}
    \textbf{Input gate:} What new information to store
    $$i_t = \sigma(W_i \cdot [h_{t-1}, x_t] + b_i)$$
    
    \vspace{0.5em}
    \textbf{Output gate:} What to output
    $$o_t = \sigma(W_o \cdot [h_{t-1}, x_t] + b_o)$$
    
    \vspace{1em}
    \textbf{Result:} Network learns what to remember and what to forget
\end{frame}

% LSTM Limitations
\begin{frame}{LSTM: Better, But Not Perfect}
    \textbf{LSTM solves vanishing gradients — mostly}
    
    \vspace{1em}
    \begin{itemize}
        \item Gradients can flow across hundreds of time steps
        \item Much better than vanilla RNN
    \end{itemize}
    
    \vspace{1em}
    \textbf{But still has limitations:}
    \begin{itemize}
        \item Doesn't remember forever — gradients still decay, just slower
        \item Sequential processing — can't parallelize across time steps
        \item Fixed-size hidden state — information bottleneck
    \end{itemize}
    
    \vspace{1em}
    \textit{These limitations motivate the search for better architectures...}
\end{frame}

% The Memory Bottleneck
\begin{frame}{The Memory Bottleneck}
    \textbf{The remaining structural problem:}
    \begin{itemize}
        \item Even LSTM compresses everything into fixed-size hidden state
        \item Real memory is addressable, structured, large
        \item A single vector is insufficient
    \end{itemize}
    
    \vspace{1em}
    \textbf{Neural Turing Machines} (DeepMind, 2014):
    \begin{itemize}
        \item RNN controller + external addressable memory bank
        \item Reads/writes via \textbf{attention mechanism}
        \item Differentiable approximation of a Turing machine
    \end{itemize}
\end{frame}

% The Paper's Prediction
\begin{frame}{The Paper's Prediction}
    \begin{quote}
        ``We expect systems that use RNNs to understand sentences or whole documents will become much better when they learn strategies for selectively attending to one part at a time.''
    \end{quote}
    
    \vspace{1em}
    \textbf{Two years later:}
    
    \vspace{0.5em}
    Vaswani et al., \textit{``Attention Is All You Need''} (2017)
    \begin{itemize}
        \item Dispenses with recurrence entirely
        \item Attention alone handles memory and context
        \item \textbf{The transformer is born}
    \end{itemize}
\end{frame}
